\scanpage{7}
\figure{\begin{tikzpicture}[x=.75cm,y=.75cm]
\foreach \t/\x/\y/\r in {U/-2/0.6/48,N/-1.65/1.03/38,I/-1.23/1.3/24,V/-.85/1.48/16,E/-.4/1.57/5,R/.1/1.6/0,S/.6/1.52/-12,I/1.05/1.36/-23,T/1.45/1.12/-33,E/1.8/.77/-45}
{\node[rotate=\r] at (\x,\y) {\t};}
\node[align=center,font=\scriptsize\bfseries] at (0,.3) {ARTS et LETTRES\\LANGUES et SCIENCES\\HUMAINES};
\foreach \t/\x/\y/\r in {P/-2/-.05/-40,A/-1.65/-.4/-32,U/-1.2/-.63/-20,L/-.75/-.8/-12,V/-.15/-.92/0,A/.3/-.9/5,L/.75/-.78/16,E/1.15/-.57/25,R/1.55/-.3/37,Y/1.85/.06/48}
{\node[rotate=\r] at (\x,\y) {\t};}
\draw (-.5,-1.65)--(.5,-1.65);
\node at (0,-2.2) {MONTPELLIER};
\end{tikzpicture}}

\begin{flushright}
MONTPELLIER, le 11 JUILLET 1984

\medskip
LE VICE-PRESIDENT DE\\
L'UNIVERSITE PAUL VALERY\\
DELEGUE AUX RELATIONS INTERNATIONALES
\end{flushright}

\uline{Service des Etudiants Etrangers}

CL/MC/N° 658

\begin{flushright}
à

\medskip
Monsieur CANO\\
Vice-Président administrateur provisoire\\
MONTPELLIER II
\end{flushright}

\bigskip
Monsieur le Président,

Dans le cadre du resserrement du jumelage entre Montpellier et Barcelone, une délégation du Conseil Municipal de Montpellier conduite par le député-maire Georges FRECHE a été reçue à Barcelone les 22 et 23 Juin dernier par Monsieur MARAGALL, maire de la ville.

Invité à accompagner cette délégation au titre de vice-président délégué aux Relations Internationales de l'Université Paul Valéry, j'ai participé à une réunion de travail présidée par M. Georges FRECHE et réunissant les vice-recteurs de l'Université Polytechnique de Catalogne, l'Université Centrale et l'Université autonome de Barcelone et l'Attaché Culturel directeur de l'Institut français de Barcelone.

Au cours de cette réunion, il a été décidé d'un commun accord de renforcer les relations pédagogiques et scientifiques entre les Universités de Barcelone et les Universités de Montpellier.

Pour ce faire, il a été convenu de recenser auprès de chaque Université les secteurs qui font déjà l'objet d'échanges entre les deux villes, ou ceux dans lesquels un travail en commun s'annoncerait fructueux. En effet, une collaboration efficace et durable entre universités repose sur une volonté et un besoin d'approfondir ensemble une activité de recherche.

A la suite de ce recensement, une rencontre des représentants des universités des deux villes sera organisée à Montpellier à l'automne 1984.

\begin{flushright}\ldots/\ldots\end{flushright}

\bigskip
Route de Mende - Boite Postale 5043 - 34032 MONTPELLIER C E D E X - Tél. (67) 63. 91. 10
\margin{6}
